\section{Main theorems}\label{sec:main-theorems}

The existence theorem concerns one lattice family, chosen before the volume sequence. We first specify its continuum conclusion; the field calculus and stress tensor established above are additional properties of the construction.

\begin{definition}[Hierarchy data]\label{def:Q0}
For a Wilson family~$\mathcal W$ with token set~$\mathcal T$, hierarchy data consist of:
\begin{itemize}[nosep]
\item complex Hilbert and vector spaces~$\mathcal H,\mathcal L$, and the relativistic data of Definitions~\ref{def:D14a}--\ref{def:D14b};
\item finite-dimensional complex normed spaces~$F_d$, linear maps~$\iota_d:F_d\to\mathcal L$, and $\pi_{\rm quo}:\mathcal T\to\mathbb F$, where $\mathbb F=\bigsqcup_dF_d$;
\item a compact metric space~$\phases$, a Borel measure~$\nu$ and $\xi_*\in\phases$;
\item functions $S_n^\zeta:\mathbb F^n\times(\R^4)^n\to\C$ for $\zeta\in\phases\sqcup\{\mathrm{mix}\}$, and $\mathsf F_n:\mathbb F^n\times\C^n\times(\R^3)^n\to\C$;
\item finite sets~$\channels{j}$, fields~$O_{j,\lambda}\in F_j$, and coefficient functions $c_{j,\lambda}:\mathbb F^I\times(\R^4)^I\times\R^4\to\C$ for every finite index set~$I$;
\item an element~$\Theta\in F_4$.
\end{itemize}
No further condition is part of these data.
\end{definition}

Put $\deg(d,A)=d$ and $\iota(d,A)=\iota_dA$. Here the spaces and maps are arbitrary: no response realization, stress interpretation, multilinearity or channel basis property is implicit. Correlations are defined at all positions; the axioms constrain their separated values.

\begin{definition}[Regular lattice approximation]\label{def:Q1}
A Wilson family is regular if its insertions are continuous and gauge invariant, each token~$v$ has a finite stencil~$\Sigma_v\subset\Z^4$ valid for every~$N,p,y$, and some~$b_*>0$ satisfies
\begin{equation*}
\forall\varepsilon>0\ \exists N_0\ \forall N\ge N_0\ \forall p:\quad
\left|\frac{\beta_{N,p}}{\log(1/a_N)}-b_*\right|<\varepsilon\,.
\end{equation*}
Only links with both ends in~$y+\Sigma_v$ modulo the torus are read. The continuum approximation is Definition~\ref{def:D09} with $K_n(v;x)=S_n^{\rm mix}(\pi_{\rm quo}(v_1),\ldots,\pi_{\rm quo}(v_n);x)$.
\end{definition}

The relativistic axioms Q2 are Definition~\ref{def:Q2}. The following conditions specify the labels, phase average and selected continuation.

\begin{definition}[Labels and phases]\label{def:Q3}
Every nonzero~$\ell\in\mathcal L$ has $\Phi(\ell,f)\ne0$ for some~$f\in\mathcal S$. Each~$\iota_d$ is injective, their images cover~$\mathcal L$, $\pi_{\rm quo}$ is onto, and some~$\Phi(\iota(a),f)$ is nonzero. The measure~$\nu$ is a probability measure of full support. Each~$S_n^\zeta$ is simultaneously permutation symmetric and smooth on~$\Conf{n}$. For every fixed separated field configuration,
\begin{equation*}
\xi\longmapsto S_n^\xi(a;x)\text{ is continuous},\qquad
S_n^{\rm mix}(a;x)=\int_{\phases}S_n^\xi(a;x)\,d\nu(\xi)\,.
\end{equation*}
\end{definition}

\begin{definition}[Selected continuation]\label{def:Q4}
The functions~$\mathsf F_n$ satisfy the ordered-time continuation condition of Definition~\ref{def:D15} and the boundary condition of Definition~\ref{def:D14f} at~$\xi_*$. In those definitions write $\mathsf F_n^{\xi_*}=\mathsf F_n$.
\end{definition}

\begin{definition}[Products]\label{def:Q5}
The single coefficient system of Definition~\ref{def:D18} satisfies both inequalities of Theorem~\ref{prop:OP} with its quantifier order: the coefficient threshold depends on $I,a,M,\eta,\varepsilon,j,\lambda$; the remainder threshold depends on $I,J,a,B,M,\eta,\delta,\varepsilon,D$. Both precede positions and states. Coefficients have no spectator, state or truncation argument.
\end{definition}

\begin{definition}[Trace test]\label{def:Q6}
For the chosen~$\Theta\in F_4$, define~$\tracetwopoint$ and~$\traceamplitude$ by the last two formulas of Definition~\ref{def:D16}. Its vacuum expectation vanishes for every Schwartz smearing, $\traceamplitude>0$, the full-radius limit of Theorem~\ref{prop:TR} holds in~$\C$, and the finite-profile exclusion of Proposition~\ref{prop:NF} holds.
\end{definition}

These conditions state the exact existence conclusion. In particular, Q6 places a condition on a degree-four field without requiring a stress-tensor interpretation.

\begin{maintheorem}[Existence]\label{thm:A}
For every gauge group of Definition~\ref{def:D01}, there is a regular \textup{(Q1)} Wilson family~$\mathcal W$ such that, for every strictly increasing~$\mathsf p:\N\to\N$, there are a strictly increasing~$\mathsf q$ and hierarchy data over~$\mathcal W$ \textup{(Q0)} satisfying \textup{Q2--Q6}, with
\begin{equation*}
\mathcal W\longrightarrow S^{\rm mix}\circ\pi_{\rm quo}
\quad\text{along }\mathsf p\circ\mathsf q
\quad\text{in the sense of Definition~\ref{def:D09}}\,.
\end{equation*}
\end{maintheorem}

The deduction below proves Theorem~\ref{thm:A} from the construction theorem. The latter retains the field calculus and quantitative estimates; subsidiary results are identified by their letter and number in Sections~\ref{sec:partII-A}--\ref{sec:partII-J}.

\begin{maintheorem}[Construction]\label{thm:B}
For every gauge group of Definition~\ref{def:D01} and~$\blockpenalty>0$, there is one regular Wilson family~$\mathcal W$, given by Definitions~\ref{def:D04} and~\ref{def:D08} along \hyperref[prop:MT]{MT}, satisfying \textup{E.1.1--E.1.3, A.1, E.2.1--E.2.1a (B1)}, and \textup{UV, E.3.1--E.3.2 (B2)}.
For every strictly increasing~$\mathsf p$, there are a strictly increasing~$\mathsf q$, the table~$S$ of \textup{IV} and hierarchy data over~$\mathcal W$ satisfying \textup{Q2--Q6}, with $\mathcal W\to S^{\rm mix}\circ\pi_{\rm quo}$ along~$\mathsf p\circ\mathsf q$ \textup{(Definition~\ref{def:D09})}. The fields~$F_d=V_d/\nullspace{d}(S)$ and tensor~\textup{ST} obey \textup{FC, F.2--F.3, \hyperref[prop:MC]{MC}, \hyperref[prop:SC]{SC}, G.2.1--G.3.3, H.1.3 (B3)}; its phases~\textup{PH} obey \textup{RC, H.2.1--H.2.4 (B4)}; its centered trace \textup{(Definition~\ref{def:D16})} obeys \textup{TR, \hyperref[prop:NF]{NF}, I.1.1--I.2.4 (B5)}; its prescription~\textup{J.1.1} obeys \textup{OP, J.1.1--J.1.6, \hyperref[prop:CT]{CT}, J.2.1--J.2.3 (B6)}.
\end{maintheorem}

The proof follows in ten steps. The degree-four position normalization and general collision-subspace estimates remain at their proof locations outside these conclusions. Restriction of the full ultraviolet limit also gives the spacing-subsequence version.

\begin{proof}[Deduction of Theorems~\ref{thm:A} and~\ref{thm:B}]
\begin{enumerate}[label=\textup{\arabic*.},itemsep=0.25em]
\item\label{step:deduction-choice} Fix~$G$ and~$\blockpenalty>0$. Appendix~\ref{sec:constants} and Sections~\ref{sec:partII-A}--\ref{sec:partII-E} supply the common choices, admissible classes, recursion and later lattice estimates. Wilson initialization \textup{D.3} starts the induction; Theorem~\ref{prop:RG} advances it; Theorem~\ref{prop:MT} and \textup{E.1.3} tune trajectories to the common endpoint. These choices precede any volume request.

\item\label{step:deduction-family} Fix the local sections, reference levels and terminal maps on those trajectories. Definition~\ref{def:D08} defines every insertion, including below its reference level. Polynomial stencil dependence gives continuity, gauge invariance and tokenwise locality. Taking~$m=N$ in Theorem~\ref{prop:MT} gives Q1 with~$b_*=\slope>0$. This fixes the regular Wilson family.

\item\label{step:deduction-uv} Proposition~\ref{prop:UB} and fixed-volume comparison with analytic tuning give Theorem~\ref{prop:UV} and \textup{E.3.1--E.3.2}, for every volume and fixed field/difference request. Thus each admissible test has a full inner limit before volume extraction.

\item\label{step:deduction-volume} Given any strictly increasing~$\mathsf p$, Theorem~\ref{prop:IV} chooses~$\mathsf q,S$, including the refinement for corrected coefficient maps. In Section~\ref{sec:partII-F}, \textup{F.1} identifies the extracted volume kernels with the full kernels by uniqueness of limits. Since~$L_{\mathsf p(\mathsf q(k))}\to\infty$, every compact support eventually fits. This proves the approximation of Definition~\ref{def:D09}, without a spacing refinement.

\item\label{step:deduction-stress} Define~$F_d=V_d/\nullspace{d}(S)$ and~$\mathcal L$ by Definition~\ref{def:D10}. Proposition~\ref{prop:FC} and \textup{F.2} give inclusions, derivatives, dimensions and the response realization. Proposition~\ref{prop:MC} and the current lemmas of Section~\ref{sec:partII-G} give the conserved line; Theorem~\ref{prop:ST} gives its normalized tensor and rotations; Proposition~\ref{prop:SC} gives the differentiated comparison. The volume refinement is unchanged.

\item\label{step:deduction-phase} Theorem~\ref{prop:PH} and \textup{H.1.1--H.1.3} of Section~\ref{sec:partII-H} construct~$\phases,\nu,S^\xi$ on the same quotient spaces. Choose~$\xi_*$ from the full-measure set with one-dimensional time-axis OS ground space. Theorem~\ref{prop:RC} reconstructs every character simultaneously; \textup{H.2.1--H.2.4} retain this selected theory and give \textup{B4}.

\item\label{step:deduction-hierarchy} Set~$\pi_{\rm quo}(d,\mathbf c)=(d,[\mathbf c])$. Representatives make this map onto. Every response lies in one~$\iota_dF_d$, and each~$\iota_d$ is injective. The unit response is nonzero since its one-point function is one; faithfulness gives a nonzero acting field. The Euclidean symmetry and smoothness conditions and Theorem~\ref{prop:PH} give Q3; reconstruction gives Q2 and Q4. These identifications use the all-spectator quotient.

\item\label{step:deduction-trace} Center the trace of this tensor in this selected phase. Theorem~\ref{prop:TR} and \textup{I.1.1--I.1.3} of Section~\ref{sec:partII-I} give its operator identity, full-radius asymptotic and unique spectral measure. Proposition~\ref{prop:NF} gives Q6's finite-profile exclusion; \textup{I.2.2--I.2.4} give the additional countertests in \textup{B5}, with~$\xi_*$ unchanged.

\item\label{step:deduction-products} Choose the prescription of \textup{J.1.1} once. Theorem~\ref{prop:OP} gives Q5; \textup{J.1.3--J.1.6} give the integrated and nested bounds in \textup{B6}. Theorem~\ref{prop:CT} and \textup{J.2.1--J.2.3} of Section~\ref{sec:partII-J} use the same underlying prescription, with separate contact projections, functionals and clearance, to give the extensions, real bases and subtraction scales. This uses its own lattice estimates.

\item\label{step:deduction-conclusion} Steps 1--9 give \textup{B1--B6} on the same family, table, quotient fields, tensor, phases and selected theory. Retain Q1--Q6 and set~$\blockpenalty=1$. Their quantifier order, lattice approximation, reconstruction, products and trace test give Theorem~\ref{thm:A}.
\end{enumerate}
\end{proof}
